% ============================================================
% I. INTRODUÇÃO
% ============================================================
\section{Introdução}
\label{sec:introducao}

As doenças cardiovasculares permanecem como a principal causa de morte no
mundo, e o eletrocardiograma (ECG) é a ferramenta não invasiva mais
amplamente utilizada para o seu diagnóstico. O volume crescente de exames,
somado à escassez de especialistas para a sua interpretação visual, tem
motivado o desenvolvimento de métodos automáticos de apoio à decisão. Em
cenários de triagem, a identificação tempestiva de condições graves ---
particularmente o infarto agudo do miocárdio (MI) e as alterações de ST/T
(STTC) --- pode impactar diretamente o desfecho do paciente, ainda que tais
exames representem fração reduzida do volume total. Métodos de aprendizado
profundo têm sido investigados como apoio à decisão clínica, alcançando
desempenho próximo ao de especialistas em tarefas
específicas~\cite{hannun2019cardiologist}.

A disponibilização da base PTB-XL~\cite{wagner2020ptbxl} como amplo
repositório público para classificação de ECG estabeleceu condições
adequadas para a comparação sistemática de algoritmos. Estudos
subsequentes~\cite{strodthoff2021benchmark} demonstraram que arquiteturas
baseadas em redes residuais convolucionais unidimensionais (ResNets 1D) e
suas variantes superam de forma consistente abordagens baseadas em
atributos manuais, estabelecendo modelos de referência sobre as
superclasses diagnósticas da PTB-XL --- NORM (ECG normal), MI, STTC,
distúrbios de condução (CD) e hipertrofia (HYP). Mecanismos complementares,
como blocos de \textit{squeeze-and-excitation}~\cite{hu2018squeeze} e
camadas de autoatenção~\cite{vaswani2017attention}, têm sido incorporados a
essas arquiteturas com ganhos consistentes e custo computacional
moderado~\cite{zhang2020seecg}.

Apesar do volume de trabalhos publicados, duas observações motivam este
estudo. Primeiro, as anotações da PTB-XL não são estritamente binárias: a
cada superclasse é associada uma confiabilidade diagnóstica atribuída pelo
cardiologista, que expressa o grau de certeza da anotação. A maior parte
dos trabalhos descarta essa informação ao binarizar os rótulos, embora ela
carregue incerteza clínica potencialmente útil. Tratar a confiabilidade
como rótulo suave (\textit{soft label}) permite que o modelo aprenda não
apenas a presença, mas o grau de suspeita de cada patologia. Segundo,
parcela relevante dos ECGs apresenta mais de uma superclasse
simultaneamente; formulações que forçam uma decisão de rótulo único
descartam informação clinicamente válida e introduzem um teto artificial de
desempenho. Uma formulação genuinamente multirrótulo, em que cada patologia
é avaliada de forma independente, é mais fiel à realidade clínica e à
estrutura do conjunto de dados.

Este trabalho propõe e avalia arquiteturas de aprendizado profundo que
estimam a confiabilidade diagnóstica contínua de cada superclasse na PTB-XL
e, a partir dela, realizam a classificação multirrótulo por limiares
otimizados. Investigam-se duas variantes de uma mesma espinha dorsal, que
funde um ramo temporal e um ramo espectral: (A) um conjunto de cinco redes
especialistas independentes, uma por superclasse; e (B) um modelo de
\textit{backbone} compartilhado com cinco saídas. As principais
contribuições são: (i) o uso explícito da confiabilidade diagnóstica como
alvo de regressão (\textit{soft labels}), com avaliação tanto da estimativa
de confiança (MAE, RMSE) quanto da classificação (AUC-ROC, AUPRC, F1
macro); (ii) uma formulação multirrótulo com limiares de decisão otimizados
por classe na validação e mantidos fixos no teste cego; e (iii) uma
comparação controlada entre a arquitetura modular de especialistas e o
\textit{backbone} compartilhado, isolando o efeito da modularidade sob o
mesmo protocolo experimental.

O restante deste artigo está organizado da seguinte forma. A
Seção~\ref{sec:relacionados} revisa trabalhos relacionados; a
Seção~\ref{sec:metodos} descreve a base de dados, as arquiteturas, a função
de perda e o protocolo experimental; a Seção~\ref{sec:resultados} apresenta
os resultados quantitativos; as Seções~\ref{sec:discussao}
e~\ref{sec:limitacoes} discutem as implicações e as limitações; e a
Seção~\ref{sec:conclusao} conclui o trabalho.
